\section{Preparación del Entorno del Taller de Desarrollo}
\begin{itemize}
    \item \codigo{mongodb}: Conexión, acceso e inserción de colecciones, campos y documentos enlazada con la base de datos situada en \textbf{MongoDB Atlas}.
    \begin{itemize}
        \item Se agrega la siguiente línea de código en \codigo{packages.json} para evitar vulnerabilidad.
        Esta configuración reemplaza la dependencia original por una versión mantenida del paquete \codigo{xlsx}.
        \begin{lstlisting}[language={JavaScript}]
"xlsx": "npm:@keep-lts/xlsx@^0.18.6"
        \end{lstlisting}    
    \end{itemize}
    \item \codigo{xlsx} para la lectura de las tablas y registros para la conversión a formato \codigo{JSON}.
    \begin{itemize}
    \item Se instalan los paquetes mencionados en el entorno local:
        El siguiente comando instala las dependencias necesarias para conectarse a MongoDB y leer archivos de Excel.
        \begin{lstlisting}
npm install xlsx mongodb
        \end{lstlisting} 
    \end{itemize}
\end{itemize}

Además, se usaron variables de entorno (\codigo{dotenv}) para proteger las credenciales y otros datos sensibles dentro del archivo \codigo{.env}.
El bloque muestra los nombres de las variables que utiliza la aplicación para establecer la conexión.
\begin{lstlisting}
DB_USER=usuario
DB_PASSWORD=password
\end{lstlisting}

Por último, se genera el script automatizado para convertir el archivo \href{https://github.com/jdcast072/template_technical-report_sena/blob/6fa21db535fab5efcd2e3cbc768ae994037af95f/others/Datafile.xlsx}{\codigo{\underline{Datafile.xlsx}}} en formato \codigo{JSON} para la creación de la base de datos e inserción de los documentos junto con sus campos por medio del lenguaje de programación de \codigo{JavaScript}.

El repositorio del script que se encuentra en el siguiente link \href{https://github.com/jdcast072/template_technical-report_sena/blob/6fa21db535fab5efcd2e3cbc768ae994037af95f/others/import-excel.js}{\codigo{\underline{import-excel.js}}}

